% SUBIR ESTE CÓDIGO COMO CAPÍTULO DE PRÁCTICA EN LA MEMORIA DE PRACTICAS
\chapter{Cuaderno de Laboratorio — Práctica 1: La Física del Aprendizaje Automático}

\section*{Objetivo}
Comprender la dinámica de la optimización mediante gradiente descendente y el impacto de los hiperparámetros en la convergencia.

\section{Tarea 1 — Importancia de la tasa de aprendizaje ($\alpha$)}

\subsection{Escenario A ($\alpha_{slow}$): convergencia lenta}
Valor elegido: $\alpha =$ \underline{\hspace{2cm}}\\
\begin{center}
\includegraphics[width=0.85\linewidth]{outputs/alpha_A_alpha_slow_a0.01.png}
\end{center}
\textbf{Análisis:} ¿Por qué baja tan lentamente? ¿Cuántas iteraciones estimas para llegar cerca del mínimo?

\vspace{2.0cm}

\subsection{Escenario B ($\alpha_{opt}$): “curva de codo”}
Valor elegido: $\alpha =$ \underline{\hspace{2cm}}\\
\begin{center}
\includegraphics[width=0.85\linewidth]{outputs/alpha_B_alpha_opt_a0.2.png}
\end{center}
\textbf{Análisis:} Describe la “curva de codo” y qué te indica sobre rapidez y estabilidad.

\vspace{2.0cm}

\subsection{Escenario C ($\alpha_{osc}$): oscilación amortiguada}
Valor elegido: $\alpha =$ \underline{\hspace{2cm}}\\
\begin{center}
\includegraphics[width=0.85\linewidth]{outputs/alpha_C_alpha_osc_a0.8.png}
\end{center}
\textbf{Análisis:} ¿Qué ocurre con $\theta$ en el espacio de búsqueda? ¿Por qué oscila y puede estabilizarse?

\vspace{2.0cm}

\subsection{Escenario D ($\alpha_{fail}$): divergencia}
Valor elegido: $\alpha =$ \underline{\hspace{2cm}}\\
\begin{center}
\includegraphics[width=0.85\linewidth]{outputs/alpha_D_alpha_fail_a2.0.png}
\end{center}
\textbf{Análisis:} ¿Por qué el error crece? Relaciónalo con “pasarse” del mínimo.

\newpage
\section{Tarea 2 — Compromiso del mini-batch}

Usa tu $\alpha_{opt}$ y compara:

\subsection{Batch completo}
\begin{center}
\includegraphics[width=0.85\linewidth]{outputs/batch_Batch_completo_b500.png}
\end{center}

\subsection{Mini-batch (32 o 16)}
\begin{center}
\includegraphics[width=0.85\linewidth]{outputs/batch_Mini_batch_32_b32.png}
\end{center}

\subsection{Estocástico puro (batch = 1)}
\begin{center}
\includegraphics[width=0.85\linewidth]{outputs/batch_Estocastico_1_b1.png}
\end{center}

\textbf{Preguntas de reflexión:}
\begin{enumerate}
    \item ¿Cuál de las tres curvas es más “ruidosa” y por qué?
    \item A nivel de tiempo de ejecución (CPU), ¿cuál ha sido más eficiente? Justifica basándote en computación vectorial.
\end{enumerate}

\section{El reto — “Ajuste de precisión” (criterio de parada)}

Implementa el criterio:
\[
|J_t - J_{t-1}| < 10^{-5}
\]
Mejor combinación: $\alpha=$ \underline{\hspace{2cm}} \quad batch=\underline{\hspace{2cm}} \quad épocas=\underline{\hspace{2cm}}.

\vspace{0.5cm}
\begin{center}
\includegraphics[width=0.85\linewidth]{outputs/reto_early_stop.png}
\end{center}

\section{Conclusiones finales}
Resume en 8--12 líneas lo aprendido sobre $\alpha$, tamaño de batch y criterio de parada.

\vspace{3.0cm}

